\chapter*{Notes}
\addcontentsline{toc}{chapter}{Notes}

\section*{Author's Afterword and Commentary}

This book is a re-imagined and carefully rewritten version of \emph{Deus Irae}, the post-apocalyptic novel begun by Philip K. Dick and completed with the collaboration of Roger Zelazny, first published in 1976. The original work---whose Latin title means \emph{God of Wrath}, a deliberate play on \emph{Dies Irae} (\emph{Day of Wrath})---stands as one of the most unusual and philosophically charged novels in Dick's body of work. Drawing on ideas first explored in Dick's short stories \emph{``The Great C''} and \emph{``Planet for Transients,''} it fused religious speculation, political trauma, and metaphysical unease in a world shattered by human decisions elevated to divine scale.

This version does not seek to replace that novel, nor to correct it. Rather, it is an act of preservation through transformation.

\subsection*{On Why This Version Exists}

The original \emph{Deus Irae} is a product of its time---brilliant, fragmented, uneven, and deeply earnest in its theological anxiety. It is also, for many contemporary readers, distant in cadence and structure, and burdened by historical context that can obscure its emotional and philosophical core. This edition was created from the conviction that the novel's central questions remain urgent, but that their presentation could be reshaped to speak more clearly to newer generations of readers.

The guiding belief behind this rewrite is simple: ideas survive best when their forms are allowed to evolve.

Accordingly, this edition modernizes language where necessary, sharpens narrative continuity, and deliberately leans into tone, rhythm, and psychological interiority---elements that resonate strongly with modern literary sensibilities. The aim has been to retain the unsettling moral pressure of the original while allowing the prose to breathe, to listen, and to unsettle in quieter, more sustained ways.

\subsection*{On the Structural Additions}

Several significant structural choices define this version:

\begin{itemize}
\item A newly written Prologue establishes the world as one where faith has already become infrastructure---regulated, monetized, and enforced---before the reader encounters the central journey. This reframes the question from \emph{``What went wrong?''} to \emph{``What did people agree to without noticing?''}

\item Three Epilogues provide closure without explanation. Rather than resolving metaphysical uncertainties, they trace consequences:
  \begin{itemize}
  \item restraint and containment of truth,
  \item kindness surviving as a minor, unofficial faith,
  \item continuation---the sense that something irreversible has begun.
  \end{itemize}

\item A sealed archival artifact reflects how institutions preserve power by deciding what must not be remembered. It reinforces the theme that history is often a managed lie, maintained as much by procedure as by belief.
\end{itemize}

These additions are not answers. They are frames, designed to let the original questions echo longer.

\subsection*{On Thematic Reorientation}

At the heart of this version is a deliberate shift in emphasis: from revelation to containment, from divine certainty to human consequence, and from apocalypse to the afterlife of belief. Truth is treated as something potentially violent, while kindness is rendered quietly subversive. Moral clarity does not arrive through exposure, but through the choices people make after certainty has failed.

The concept of the ``membrane,'' introduced elliptically and never explained, functions as a metaphor for this transition: a world beginning to listen back, not through miracles, but through attention, clarity, and discomfort.

\subsection*{On Title, Legacy, and Intent}

While retaining the conceptual inheritance of \emph{Deus Irae}, this version reframes the story away from wrath as spectacle and toward wrath as residue---something that lingers in systems, habits, and stories long after its source is gone. The God of Wrath here is not defeated; he is archived. That distinction matters.

This edition exists to honor a masterpiece by refusing to fossilize it. Philip K. Dick's work has always thrived on instability, reinterpretation, and the suspicion that reality itself is provisional. To present \emph{Deus Irae} as fixed and untouchable would be to betray its spirit.

This book is offered instead as a bridge: between generations of readers, between faith and doubt, between the need to know and the need to live with uncertainty.

If it leads readers back to the original novel, it has succeeded. If it leads others to encounter these ideas for the first time, it has done its work.

A face can become law.  
A story can become a prison.  
But a question---properly asked---can still open a door.
